% ════════════════════════════════════════════════════════════════
\subsection{Termin poprawkowy}
% ════════════════════════════════════════════════════════════════

\begin{zadanie}[Zadanie 1 --- niemalejące długości najkrótszych ścieżek powiększających (15 p.)]
	$P$ --- najkrótsza ścieżka powiększająca dla $M$, $M'=M\div P$, $P'$ ---
	najkrótsza ścieżka powiększająca dla $M'$. Udowodnij $|P|\le|P'|$ oraz że
	wspólny wierzchołek wymusza $|P|<|P'|$.
\end{zadanie}

(Lemat z~algorytmu Hopcrofta--Karpa.) Mierzymy $|P|$ liczbą krawędzi.

\paragraph{Nierówność $|P|\le|P'|$.}
Niech $M''=M'\div P'$; wtedy $|M''|=|M|+2$. Rozważmy $D=M\div M''$. Jako
symetryczna różnica dwóch skojarzeń, $D$ ma stopnie $\le 2$, więc rozpada się na
ścieżki i~cykle naprzemienne względem $M$. Skoro $|M''|=|M|+2$, w~$D$ jest o~dwie
krawędzie $M''$ więcej niż $M$, zatem $D$ zawiera co najmniej dwie
wierzchołkowo-rozłączne ścieżki powiększające $Q_1,Q_2$ względem~$M$. Każda
z~nich ma długość $\ge|P|$ (bo $P$ jest najkrótsza), więc
\[
	|D|\ \ge\ |Q_1|+|Q_2|\ \ge\ 2|P|.
\]
Z~drugiej strony
\[
	D=M\div M''=(M\div M')\div(M'\div M'')=P\div P',
\]
a~dla symetrycznej różnicy zbiorów krawędzi
$|P\div P'|=|P|+|P'|-2\,c$, gdzie $c=|E(P)\cap E(P')|$ to liczba wspólnych
krawędzi. Łącząc:
\[
	|P|+|P'|-2c\ \ge\ 2|P|
	\quad\Longrightarrow\quad
	|P'|\ \ge\ |P|+2c\ \ge\ |P|.
\]

\paragraph{Wspólny wierzchołek $\Rightarrow$ ostra nierówność.}
Pokażemy kontrapozycję: jeśli $|P'|=|P|$, to $P$ i~$P'$ są wierzchołkowo
rozłączne. Z~powyższego $|P'|=|P|$ wymusza $c=0$ (brak wspólnych krawędzi),
więc $D=P\div P'$ jest rozłączną krawędziowo sumą $P$ i~$P'$. Ale $D$ ma stopnie
$\le 2$:
\begin{itemize}[nosep]
	\item gdyby $P$ i~$P'$ miały wspólny wierzchołek wewnętrzny $v$, jego stopień
	      w~$D$ wynosiłby $2+2=4>2$ --- sprzeczność;
	\item wspólnym wierzchołkiem nie może też być koniec: końce $P$ są wolne
	      względem $M$, lecz \uemph{skojarzone} względem $M'=M\div P$, więc nie są
	      wolnymi końcami $P'$ (a~końce $P'$ muszą być wolne względem $M'$); końce
	      $P'$ z~kolei mają w~$D$ stopień $1$, a~wierzchołek wewnętrzny $P$ ma
	      stopień~$2$, co znów dałoby stopień~$3$.
\end{itemize}
Zatem $P,P'$ są wierzchołkowo rozłączne. Równoważnie: wspólny wierzchołek
$\Rightarrow |P'|\ne|P|$, a~ponieważ $|P'|\ge|P|$, otrzymujemy $|P|<|P'|$. $\qed$

\clearpage
\begin{zadanie}[Zadanie 2 --- matroid podzbiorów krawędzi o~ograniczonym stopniu (15 p.)]
	$G$ spójny dwudzielny, $S\in\mathcal{F}\iff$ każdy wierzchołek incydentny
	z~$\le k$ krawędziami z~$S$. Czy $(E,\mathcal{F})$ zawsze jest matroidem?
\end{zadanie}

\paragraph{Nie --- kontrprzykład.}
Dla $k=1$ rodzina $\mathcal{F}$ to rodzina skojarzeń w~$G$. Weźmy ścieżkę
$P_4$ na czterech wierzchołkach z~krawędziami $e_1,e_2,e_3$ (kolejno);
jest to graf spójny i~dwudzielny. Rozważmy
\[
	I_1=\{e_2\}\in\mathcal{F},\qquad I_2=\{e_1,e_3\}\in\mathcal{F},
	\qquad |I_1|<|I_2|.
\]
Aby zachować własność wymiany, powinniśmy móc dodać do $I_1$ pewną krawędź
z~$I_2\setminus I_1$. Jednak $e_1$ i~$e_2$ mają wspólny wierzchołek (stopień $2>1$),
tak samo $e_2$ i~$e_3$. Zatem ani $\{e_2,e_1\}$, ani $\{e_2,e_3\}$ nie należą do
$\mathcal{F}$ --- własność wymiany jest złamana. Para $(E,\mathcal{F})$ nie jest
więc matroidem. $\qed$

\clearpage
\begin{zadanie}[Zadanie 3 --- mnożenie dużych liczb metodą Tooma-Cooka (15 p.)]
	Wykorzystując dane wzory (3 części, mnożenia $x_2y_2$, $x_0y_0$, $A$, $B$, $C$),
	zaprojektuj algorytm mnożenia dużych liczb i~oszacuj jego czas.
\end{zadanie}

\paragraph{Algorytm (Toom--Cook 3-way).}
Liczbę $n$-cyfrową zapisujemy w~bazie $10^{k}$, $k=\lceil n/3\rceil$, jako
$x_2 10^{2k}+x_1 10^{k}+x_0$ (i~analogicznie dla drugiej liczby); każda część
$x_i,y_i$ ma $\approx n/3$ cyfr. Liczymy pięć iloczynów
\[
	x_2y_2,\quad x_0y_0,\quad A,\quad B,\quad C,
\]
gdzie $A,B,C$ to iloczyny sum/różnic części z~małymi stałymi współczynnikami
($2,4$); mnożniki te mają nadal $\BO(n/3)$ cyfr, więc to $5$ mnożeń liczb
rozmiaru $\approx n/3$ (rekurencyjnie). Z~tych pięciu iloczynów składamy wynik
danymi wzorami: pozostałe operacje to dodawania, odejmowania, przesunięcia
(mnożenia przez $10^{jk}$) oraz dzielenia przez \uemph{stałe} $2$ i~$6$ ---
wszystko w~czasie $\BO(n)$.

\paragraph{Złożoność.}
\[
	T(n)=5\,T\!\Big(\tfrac n3\Big)+\BO(n).
\]
Ponieważ $\log_3 5\approx 1{,}465>1$, dominuje człon rekurencyjny:
\[
	T(n)=\Theta\!\big(n^{\log_3 5}\big)\approx\Theta\!\big(n^{1{,}465}\big),
\]
czyli asymptotycznie szybciej niż Karacuba ($n^{\log_2 3}\approx n^{1{,}585}$).
$\qed$

\clearpage
\begin{zadanie}[Zadanie 4 --- najkrótszy cykl Hamiltona (15 p.)]
	Algorytm znajdujący długość najkrótszego cyklu Hamiltona w~grafie ważonym
	o~$n$ wierzchołkach w~czasie $\BO(2^n)$ (z~dokładnością do~wielomianu).
\end{zadanie}

\paragraph{Programowanie dynamiczne Helda--Karpa.}
Ustalamy wierzchołek startowy $1$. Dla $S\subseteq V$ z~$1\in S$ oraz $v\in S$
niech $D[S][v]$ będzie minimalną wagą ścieżki, która zaczyna się w~$1$, odwiedza
\uemph{dokładnie} zbiór $S$ i~kończy się w~$v$. Wtedy
\[
	D[\{1\}][1]=0,\qquad
	D[S][v]=\min_{u\in S\setminus\{v\}}\Big(D\big[S\setminus\{v\}\big][u]+w(u,v)\Big).
\]
Długość najkrótszego cyklu Hamiltona to
\[
	\min_{v\neq 1}\Big(D[V][v]+w(v,1)\Big).
\]

\paragraph{Poprawność.}
$D[S][v]$ poprawnie liczy optymalną ścieżkę przez $S$ kończącą się w~$v$:
przedostatni wierzchołek $u$ jednoznacznie wyznacza prefiks odwiedzający
$S\setminus\{v\}$ i~kończący się w~$u$, optymalizowany rekurencyjnie
(podstruktura optymalna). Dodanie krawędzi powrotnej $w(v,1)$ domyka cykl
Hamiltona; minimalizacja po~$v$ daje optimum.

\paragraph{Złożoność.}
Stanów $D[S][v]$ jest $\BO(2^n n)$, a~każdy liczymy w~czasie $\BO(n)$:
\[
	\BO(2^n n^2)=\BO^*(2^n). \qquad\qed
\]

\clearpage
\begin{zadanie}[Zadanie 5 --- 3-aproksymacyjny zbiór niezależny w~grafach dysków jednostkowych (15 p.)]
	Wielomianowy algorytm $3$-aproksymacyjny dla najliczniejszego zbioru
	niezależnego w~grafach dysków jednostkowych.
\end{zadanie}

Korzystamy z~tego samego faktu geometrycznego, co przy kolorowaniu
(podrozdz.~\ref{sol:udg-coloring}): zbiór punktów z~prawej półpłaszczyzny koła
o~promieniu $1$ wokół najbardziej lewego punktu $p$ rozkłada się na $3$ wycinki
$60^\circ$, z~których każdy jest kliką.

\paragraph{Algorytm (zachłanne ,,obieranie'' od~lewej).}
\begin{enumerate}[nosep]
	\item $\mathcal{I}\leftarrow\emptyset$.
	\item Dopóki zostały punkty: wybierz najbardziej lewy (najmniejsze $x$) punkt
	      $p$, dodaj $p$ do $\mathcal{I}$ i~usuń $p$ oraz wszystkich jego sąsiadów
	      $N(p)$.
	\item Zwróć $\mathcal{I}$.
\end{enumerate}
$\mathcal{I}$ jest niezależny: dodając $p$ usuwamy wszystkich jego sąsiadów, więc
żadne dwa wybrane punkty nie sąsiadują.

\paragraph{Współczynnik aproksymacji.}
Niech $p_1,p_2,\dots,p_t$ będą kolejno wybranymi punktami; usuwane domknięte
sąsiedztwa $N[p_i]=\{p_i\}\cup N(p_i)$ są parami rozłączne i~pokrywają wszystkie
wierzchołki. Niech $\mathrm{OPT}$ będzie najliczniejszym zbiorem niezależnym.
Dla każdego $i$ punkty $N(p_i)$ leżą (w~chwili wyboru, gdy $p_i$ jest najbardziej
lewy) w~prawej półpłaszczyźnie koła, więc rozkładają się na $3$ kliki. Zbiór
niezależny zawiera co najwyżej jeden punkt z~każdej kliki, więc
\[
	|\mathrm{OPT}\cap N[p_i]|\ \le\ 3.
\]
(Jeśli $p_i\in\mathrm{OPT}$, to żaden jego sąsiad nie jest w~$\mathrm{OPT}$, więc
wkład wynosi $1$.) Ponieważ sąsiedztwa $N[p_i]$ pokrywają wszystkie wierzchołki
i~są rozłączne,
\[
	|\mathrm{OPT}|=\sum_{i=1}^t|\mathrm{OPT}\cap N[p_i]|\le 3t=3|\mathcal{I}|,
\]
czyli $|\mathcal{I}|\ge|\mathrm{OPT}|/3$. Algorytm jest $3$-aproksymacyjny
i~działa w~czasie wielomianowym. $\qed$
